\begin{center}
\textsc{Gauss} to \textsc{Olbers}.\quad Göttingen, 1~March 1827.
\end{center}

\par\noindent\hspace*{\parindent}\mbox{.\ .\ .\ .\ .\ .}\quad I have finished my paper, or perhaps more properly my first paper, on curved surfaces, but I shall not yet submit it to the Society, since in any event no volume
\par\smallskip\noindent\hspace*{0.03\linewidth}{\footnotesize\textbf{IX.}}\hfill{\footnotesize48}\hspace*{0.14\linewidth}\par
\clearpage
\noindent\makebox[0pt][l]{\fontsize{13}{16}\selectfont 378}\hfill{\small\textbf{CORRESPONDENCE WITH OLBERS.}}\hfill\null\par\medskip
\noindent will come out by the Easter fair.
Printing has not yet begun on the two papers I submitted in 1825 and 1826, on biquadratic residues and Suppl.\ theor.\ combin.\ observ.
The paper on curved surfaces really contains only a couple of propositions for \textls[120]{direct} use in my future work on the survey, namely 1)~what is required to calculate the excess of the sum of the 3 angles over \(180^{0}\) in a triangle whose sides are shortest lines on a \textls[120]{non}-spherical surface, and 2)~how in this case the excess must be distributed \textls[120]{unequally} among the three angles so that their sines become proportional to the opposite sides.
In practice this is quite unimportant, since even in the largest triangles that can be measured on Earth this inequality in the distribution is imperceptible. The dignity of science nevertheless demands that the nature of this inequality be clearly understood.
And so here, as on many other occasions, one may indeed exclaim Tantae molis erat! at the effort of getting that far. ---
More important than the solution of these 2 problems, however, is that the paper establishes several general principles from which the solution of a great many important problems can later be derived in a more specialized investigation.

\par\noindent\hspace*{\parindent}I return once more to the subject mentioned at the beginning of this letter.
I have based all my auxiliary tables and calculations on \textsc{Walbeck}'s flattening of \(\frac{1}{302{,}78}\), but I believe that all the arc measurements made so far, if their data were taken into account \textls[120]{in full}, would show that \textsc{Sabine}'s flattening of \(\scriptstyle\frac{1}{2\mkern5mu 8\mkern5mu 9}\) could be reconciled with them almost as well, and it seems to me that all existing arc measurements still cover far too small an extent to confine the flattening within narrower limits.
What is certain is that the Earth is an irregular body. The altitudes of the pole at different places (leaving errors of observation entirely aside) always vary by several arcseconds (perhaps quite a few arcseconds here and there) around the values one \textls[120]{calculates} on the assumption of some regular ellipsoid, and likewise the actual pendulum lengths vary around the calculated ones (the \textls[120]{mean} variation in the pendulum lengths is perhaps \(\scriptstyle\frac{1}{3\mkern7mu 0}\) of an English line).
But for that very reason arc measurements of small extent can contribute little to our knowledge of the mean figure of the Earth. In particular, I consider the Lapland
\par\clearpage
\noindent\null\hfill{\small\textbf{ON THE HANOVER TRIANGULATION.}}\hfill\makebox[0pt][r]{\fontsize{13}{16}\selectfont \textbf{379}}\par\medskip
\noindent arc far too small.
This makes it all the more important, I think, that all precise triangulations in Europe should gradually be linked into a single whole.
I now have hopes of receiving the Bavarian triangles.
Once all the European observatories from Åbo to Palermo and from Nicolajef to Dublin are linked by triangles, so that their relative positions are determined with all the accuracy that the finest geodetic measurements can afford, it will be possible, i.e.\ it will be possible for our descendants, to judge everything with much greater \mbox{certainty.\ .\ .\ .\ .\ .\ .}
